% STATUS: verified
\chapter{Architecture}\label{ch:arch}

This chapter describes the system as a set of running parts: the eight containers and what each
owns, the two networks and the single way in and the single way out, the four rules that decide
which part may do what, and the two internal doors between the API and the assistant service.
It then follows four requests through the parts --- a chat turn, a report, a work product and a
refresh of the legal corpus --- at the level of who calls whom. Background jobs are described
in full here, because every long operation of the system is one: the queue, the worker, the
scheduler and the daily maintenance. The chapter closes with the technology stack, the two
modes of the compose file, and the architectural decisions with their reasons.

\section{The system and its containers}\label{sec:arch:containers}

The \product{} helps a compliance officer establish which Lithuanian and EU acts apply to a
financial-sector company, extracts the obligations those acts impose, compares them with the
company's own documents and writes a report of the gaps. From the same documents it drafts the
work products of a compliance programme for a person to edit and approve. The user works in a
chat; structured pages exist for review and administration. All of it runs on one host as one
Docker Compose project, described in \file{docker-compose.yml}: eight containers
(\cref{tab:arch:containers}), of which three are built from the repository's own code --- the
API, the assistant service and the console --- and one, the database, from a stock image with
an extension compiled in.

\begin{table}[H]
\caption{The eight containers of \file{docker-compose.yml} and what each owns}\label{tab:arch:containers}
\begin{tabularx}{\textwidth}{@{}L{36mm}L{33mm}Y@{}}
\toprule
\hd{Container} & \hd{Image} & \hd{What it owns} \\
\midrule
\code{legalops-endpoint} & \code{caddy:2.11-alpine} & The single ingress. It publishes host port 80 and routes each request by its path to the console, the API, the assistant service or the database tool; it sets the security headers and writes a JSON access log. \\
\code{legalops-exit} & \code{vuknf/caddy-l4:2.11} & The single egress: a proxy on port 443 that forwards a TLS connection only when it names one of four allowed hosts. \\
\code{legalops-vite} & built from \file{vite/} & The console under \code{/console} and the public site at the root: one Vite project with two applications. It holds no data and no secret. \\
\code{legalops-django} & built from \file{django/} & The API (Django): accounts and sessions, companies and their documents, the legal corpus, the register of obligations, reports, work products, prompts, chat history, usage and the audit log; the uploaded files on disk; the internal door \code{/internal/assistant/*}. The worker's loop runs inside this container. \\
\code{legalops-worker} & \code{docker:27-cli} & Nothing of its own. It is a Docker command-line client that starts the job loop inside \code{legalops-django} and stays alive as long as that loop does. \\
\code{legalops-assistant} & built from \file{assistant/} & The assistant service: every call to a model --- the streaming chat with its tools, the generations and embeddings the API asks for, and the client of external MCP servers. It owns the gateway's key and no data. \\
\code{legalops-postgres} & built from \file{postgres/} & PostgreSQL with the \code{pgvector} extension: all data, the vector indexes and the full-text indexes. An empty data directory loads \file{postgres/bootstrap.sql.gz} on the first start. \\
\code{legalops-dbgate} & \code{dbgate/dbgate:7.3.3-alpine} & An administration tool: a browser for the database, reachable under \code{/dbgate/} by administrators only. \\
\bottomrule
\end{tabularx}
\end{table}

Every container has a read-only root file system; what must be written goes to a
\code{tmpfs} mount or to one of the two host folders under \file{_DATA/} (the database files
and the uploads). The endpoint, the exit, the API and the database run as uid~1000. The
worker and the database tool run as root, and for the console's and the assistant service's
containers the \code{user} line of the compose file is commented out, so they run as their
image's default user. The compose file sets no resource limits, no health checks and no start
order: the services start together, and the endpoint answers 502 until the service behind a
path is up. The security consequences of these choices are the subject of \cref{ch:security};
installation and updates are in \cref{ch:ops}.

\paragraph{Inside the API.} The Django project uses none of Django's bundled applications: no
admin site, no \code{django.contrib.auth}, no sessions framework and no templates of its own
(\file{settings.py}). It is eleven applications of the project's own (\cref{tab:arch:apps}),
two pieces of middleware --- \code{RequestGuardMiddleware} (\file{common/middleware.py}),
which refuses a NUL character in a path or query and turns a value the database cannot take,
one out of range or malformed, into
a 400 answer, and Django's \code{CommonMiddleware} --- and one route table for the whole
service, \file{urls.py}. Every answer is JSON, errors included. Times are aware UTC values
throughout. The routes themselves are listed in \cref{ch:api} and the tables in \cref{ch:data}.

\begin{table}[H]
\caption{The applications of the API, as \file{settings.py} installs them}\label{tab:arch:apps}
\begin{tabularx}{\textwidth}{@{}L{23mm}Y@{\hspace{5mm}}L{23mm}Y@{}}
\toprule
\hd{Application} & \hd{Holds} & \hd{Application} & \hd{Holds} \\
\midrule
\code{users} & Accounts, sessions, sign-in & \code{programme} & Work products \\
\code{audit} & The audit log & \code{chat} & Threads and messages \\
\code{jobs} & The queue and the worker & \code{prompts} & Prompts and their versions \\
\code{companies} & Companies and their documents & \code{ops} & Health probe, maintenance, rate-limit windows \\
\code{corpus} & Acts, versions, units, obligations & \code{assistant} & The client of the assistant service, usage, prices, MCP servers \\
\code{analysis} & Reports and the action register & & \\
\bottomrule
\end{tabularx}
\end{table}

\paragraph{Inside the assistant service.} The service is one Node process: an Express
application (\file{app.js}) served by \file{server.js}. Its folders follow its three jobs:
\file{routes/} holds the public chat route and the two internal route groups, \file{llm/} the
provider registry, the chat tools, the stream and the MCP client, and \file{services/} the
clients of the API and the small in-memory registries (who is signed in, the turn limit, the
live turns). It keeps nothing on disk.

\section{Networks, the way in and the way out}\label{sec:arch:network}

The containers sit on two Docker networks, and the difference between them is the main
security property of the deployment (\cref{fig:arch:containers}).

\begin{description}
  \item[The isolated network] \code{isolated} is a bridge declared \texttt{internal: true} with
    the subnet 172.28.10.0/24. Docker gives such a network no route to the outside. Every
    service except the worker is attached to it, and it is the only network the API, the
    assistant service, the database and the database tool are attached to in production
    mode.
  \item[The external network] \code{external} is an ordinary Docker network that the compose
    file does not create itself (it is declared \texttt{external: true};
    \file{runUpdateThisStack.sh} creates it when the host has none of that name). Only the
    endpoint and the exit join it in
    production mode. The endpoint is the one container with a published port.
  \item[No network] The worker has \texttt{network\_mode: none}. It talks to nothing but the
    Docker socket mounted into it.
\end{description}

\begin{figure}[H]
\centering
\begin{tikzpicture}[x=1mm, y=1mm]
  \begin{scope}[on background layer]
    \draw[lane] (0,-14) rectangle (150,-86);
  \end{scope}
  \node[lanelabel, anchor=north west] at (1.5,-15.5) {isolated network --- internal, no route out};

  \node[ext] (browser) at (75,6) {browser, through the hosting's TLS proxy};
  \node[svc, text width=24mm] (endpoint) at (75,-14) {\textbf{endpoint}\\ Caddy, port 80};
  \node[svc, text width=22mm] (vite) at (16,-40) {\textbf{vite}\\ public site, console};
  \node[svc, text width=22mm] (dbgate) at (134,-40) {\textbf{dbgate}\\ database tool};
  \node[svc, text width=24mm] (assistant) at (40,-63) {\textbf{assistant}\\ every model call};
  \node[svc, text width=24mm] (django) at (98,-63) {\textbf{django}\\ the API and the worker's loop};
  \node[store] (postgres) at (136,-63) {\textbf{postgres}};
  \node[svc, text width=22mm] (exit) at (62,-86) {\textbf{exit}\\ SNI allowlist};
  \node[svc, text width=22mm] (worker) at (136,-101) {\textbf{worker}\\ Docker CLI, no network};

  \node[ext] (gw) at (13,-109) {ai.knf.vu.lt};
  \node[ext] (cellar) at (48,-109) {publications.europa.eu};
  \node[ext] (etar) at (83,-109) {www.e-tar.lt};
  \node[ext] (lb) at (106,-109) {www.lb.lt};

  \draw[flow] (browser) -- (endpoint);
  \draw[flow] (endpoint) -- node[lbl, pos=.55] {everything else} (vite);
  \draw[flow] (endpoint) -- node[lbl, pos=.5] {\code{/api/assistant/*}} (assistant);
  \draw[flow] (endpoint) -- node[lbl, pos=.5] {\code{/api/*}} (django);
  \draw[flow] (endpoint) -- node[lbl, pos=.55] {\code{/dbgate/*}} (dbgate);

  \draw[flow] ([yshift=2.2mm]django.west) -- node[lbl, above] {\code{/internal/llm/*}} ([yshift=2.2mm]assistant.east);
  \draw[flow] ([yshift=-2.2mm]assistant.east) -- node[lbl, below, align=center] {\code{/api/*} as the user,\\ \code{/internal/assistant/*}} ([yshift=-2.2mm]django.west);
  \draw[flow] (django) -- (postgres);
  \draw[flow] (dbgate) -- (postgres);

  \draw[flow] (assistant.south) -- (exit.north west);
  \draw[flow] (django.south west) -- (exit.north east);
  \draw[flow] (exit.south west) -- (gw.north);
  \draw[flow] (exit.south) -- (cellar.north);
  \draw[flow] (exit.south east) -- (etar.north);
  \draw[flow] (exit.east) -- (lb.north);

  \draw[opt] (worker.north west) -- node[lbl, pos=.5] {\texttt{docker exec}} (django.south east);

  \node[note, anchor=west] at ([xshift=1mm]endpoint.east) {on both networks;\\ the only published port};
  \node[note, anchor=east] at ([xshift=-1mm]exit.west) {on both networks,\\ at 172.28.10.30};
\end{tikzpicture}
\caption{The containers (their names without the prefix \code{legalops-}), the isolated
  network, and which service may reach which. An arrow points from the caller to the service
  called. Requests enter only through the endpoint and leave only through the exit, which
  forwards to four named hosts. The dotted arrow is not a network path: the worker starts the
  job loop inside the API's container through the Docker socket.}\label{fig:arch:containers}
\end{figure}

\subsection{The single ingress}

The endpoint is a Caddy reverse proxy configured in \file{endpoint/Caddyfile}. It listens on
port 80 in plain HTTP; TLS ends at the hosting's proxy in front of the host. It first refuses
a request whose \code{Origin} header names another site, with a bare 403, and adds a content
security policy
and a transport security header to every other answer (\cref{ch:security}), then routes by
path, first match wins
(\cref{tab:arch:ingress}).

\begin{table}[H]
\caption{How the endpoint routes a request, in the order of \file{endpoint/Caddyfile}}\label{tab:arch:ingress}
\begin{tabularx}{\textwidth}{@{}L{34mm}L{45mm}Y@{}}
\toprule
\hd{Path} & \hd{Goes to} & \hd{Remark} \\
\midrule
\code{/dbgate/*} & \code{legalops-dbgate:3000} & Only after the API has confirmed, on \route{GET}{/api/auth/checkauth/admin}, that the caller is a signed-in administrator (Caddy's \code{forward_auth}). \\
\code{/api/assistant/*} & \code{legalops-assistant:3000} & The chat stream and its stop and status routes. \\
\code{/api/*} & \code{legalops-django:8000} & The whole public API. \\
everything else & \code{legalops-vite:80} & The public site and the console. \\
\bottomrule
\end{tabularx}
\end{table}

Two things follow from this table. The assistant service has exactly one public route group,
the chat; everything else it serves is internal. And no path beginning with \code{/internal/}
is forwarded to the API or to the assistant service: such a request falls through to the last
row and is answered by the console's container, which serves only the public site and the
console.

\subsection{The single egress}

A container on the isolated network cannot open a connection to the internet. The exit is
the one exception that is allowed, and it is narrow. It is a layer-4 proxy
(\file{exit/Caddyfile}) that listens on port 443 only, reads the server name a TLS client
announces when it opens a connection (the SNI field) and, when that name is on its list,
forwards the connection to the same name on port 443. It does not decrypt anything. A
connection that names any other host matches no route and is not forwarded.

A service reaches an allowed host through two entries that must both exist: a route in
\file{exit/Caddyfile}, and an \code{extra_hosts} line on the consuming service in
\file{docker-compose.yml} that pins the host's name to the exit's fixed address,
172.28.10.30. The service then connects to what it believes is the host and reaches the exit.
\Cref{tab:arch:egress} lists the four hosts.

\begin{table}[H]
\caption{The egress allowlist: the hosts of \file{exit/Caddyfile} and the services that have
  them pinned}\label{tab:arch:egress}
\begin{tabularx}{\textwidth}{@{}L{38mm}L{30mm}Y@{}}
\toprule
\hd{Host} & \hd{Pinned for} & \hd{Used for} \\
\midrule
\code{ai.knf.vu.lt} & assistant service, API & The gateway: every model call, made by the assistant service. The compose file pins the name for the API's container too, but no code of the API calls it. \\
\code{publications.europa.eu} & API & CELLAR, the Publications Office's machine interface: the EU acts. \\
\code{www.e-tar.lt} & API & The Lithuanian register of legal acts. \\
\code{www.lb.lt} & API & The Bank of Lithuania's pages and PDF files. \\
\bottomrule
\end{tabularx}
\end{table}

The three corpus hosts are pinned for the API's container because the worker's loop, which
fetches the acts, runs there. The fetcher adds its own rule on top of the network's
(\file{corpus/fetchers.py}): HTTPS on port 443 only, checked again on every redirect, so that
a source cannot point the worker at one of the stack's own services; which hosts are
reachable at all is left to the allowlist.

The reason for the arrangement is that the two services that hold the system's data and its
gateway key also process text they did not write --- uploaded documents, fetched pages,
the answers of a model, the tools of external MCP servers. If one of them were made to open a
connection of an attacker's choosing, the isolated network gives that connection nowhere to
go except the four hosts. The price is paid in operations: a new corpus source or an external
MCP server works only after its host is added in both places (\cref{ch:ops}). EUR-Lex and
e-seimas are deliberately absent, because they refuse unattended clients (\cref{ch:corpus}).

\section{Four design rules}\label{sec:arch:rules}

Four rules fix the division of work between the API and the assistant service. Each exists
for a reason and each has a consequence that can be checked in the code.

\begin{description}
  \item[The assistant service owns no data.]
    Prompts, chat threads, usage records and the configurations of MCP servers live in the
    API's database. The assistant service has no database connection --- its container is
    given no database address --- and reads what it needs through the API. \emph{Reason:}
    one copy of every fact, behind one set of access rules, in one place.
    \emph{Consequence:} the service can be restarted or replaced without losing anything
    but the chat turns in flight. What it keeps in memory --- the turn limiter's windows, the
    cache of signed-in callers, the registry of live turns --- is rebuilt from nothing, and
    the code notes that it therefore runs as a single instance on purpose
    (\file{services/runs.js}, \file{services/ratelimit.js}).

  \item[The API never calls a model provider.]
    The worker and the web process ask the assistant service for a generation
    (\route{POST}{/internal/llm/generate}) or for embeddings
    (\route{POST}{/internal/llm/embed}); the one client is \file{assistant/gateway.py} in the
    API. The API makes outbound HTTP requests from exactly two modules: this client, and
    the corpus fetcher. \emph{Reason:} the providers and their credentials are configured
    once, in one container. \emph{Consequence:} the gateway's key is in the environment of
    the assistant service alone, a model is named everywhere by one string of the form
    \code{provider:model} stored on the prompt version, and every model call of the system
    passes one piece of code that can time it, retry it and count its tokens
    (\cref{ch:models}).

  \item[Chat tools act through the public API with the user's own token.]
    When the model in a chat calls a tool --- file a company profile, read a report, start
    an analysis --- the assistant service sends an ordinary request to \code{/api/*} carrying
    the bearer token of the user who is chatting (\code{userClient} in
    \file{services/django.js}). \emph{Reason:} the rule of who may see what --- a user sees
    their own rows, an administrator all rows except other users' chat threads --- is
    enforced once, in the API's views, and a tool cannot step around it. \emph{Consequence:}
    the chat can do nothing its user could not do in the console, every change it makes is
    audited under that user's name, and a tool has no route to accounts, prompts or MCP
    settings, and none that changes the corpus: the chat may search and read the acts and
    the register of obligations, never edit them (\cref{ch:chat}). There are two deliberate exceptions, both through the
    internal door described below: the assistant's own replies, and the title a model gives
    a conversation.

  \item[Every model call is logged.]
    A call leaves a row in \code{llm_calls} with its model, its token counts, its cost in US
    dollars from the gateway's price list, the prompt version and the outcome. The API
    writes the row itself after each generation and each embedding batch it asked for
    (\code{record_call} in \file{assistant/usage.py}); the assistant service reports each
    chat turn, with its transcript, through the internal door. \emph{Reason:} cost and
    accountability --- a budget can only be enforced, and an answer only be traced, from a
    complete record. \emph{Consequence:} the daily budget of a user, the usage page and the
    cost of every job step are all computed from this one table. Two limits of the rule
    should be stated: the embedding of a search query (a few tokens) is not recorded, and a
    generation that fails raises before the row is written, so the rows of job steps are
    those of calls that answered.
\end{description}

\section{The two internal doors}\label{sec:arch:doors}

The API and the assistant service call each other, and neither call is open to anyone else.
Each side has an \term{internal door}: a group of routes under \code{/internal/} that exists
only for the other service (\cref{tab:arch:doors}).

\begin{xltabular}{\textwidth}{@{}L{62mm}Y@{}}
\caption{What goes through each internal door}\label{tab:arch:doors}\\
\toprule
\hd{Route} & \hd{What goes through} \\
\midrule
\endfirsthead
\toprule
\hd{Route} & \hd{What goes through} \\
\midrule
\endhead
\bottomrule
\endlastfoot
\multicolumn{2}{@{}L{\textwidth}@{}}{\hd{The API's door} --- called by the assistant service (\code{internalFetch} in \file{services/django.js}), served by \file{assistant/api/internal_views.py}} \\
\route{GET}{/internal/assistant/prompt} & The active version of one prompt with its text and model. When a prompt has none the answer is 404 and the chat answers 503 rather than run on a text nobody can see. \\
\route{GET}{/internal/assistant/budget} & What a user may still spend today; asked before each chat turn. \\
\route{POST}{/internal/assistant/turn-log} & One \code{llm_calls} row for a chat turn: tokens, the tools called, the outcome and the transcript. \\
\route{POST}{/internal/assistant/threads/{id}/messages} & The assistant's reply --- text, tool calls and tool results --- stored when the turn ends. Only this door may write a message of the role \code{assistant}. \\
\route{POST}{/internal/assistant/threads/{id}/title} & The short title a model wrote for a new conversation; applied only while the thread still carries its automatic title. \\
\route{GET}{/internal/assistant/mcp-servers} & The enabled MCP servers with the header values that authenticate to them --- the only place those values leave the API. \\
\midrule
\multicolumn{2}{@{}L{\textwidth}@{}}{\hd{The assistant service's door} --- called by the API (\file{assistant/gateway.py}), served by \file{routes/internal.js} and \file{routes/internalMcp.js}} \\
\route{POST}{/internal/llm/generate} & One generation: a prompt, optional instructions and shared data, a model, and optionally a JSON schema the answer must fit. Answers the text or the object, with the usage. \\
\route{POST}{/internal/llm/embed} & A batch of at most 256 texts (\code{MAX_EMBED_TEXTS}) to embed. \\
\route{GET}{/internal/llm/models} & The configured providers and the gateway's models with their prices. \\
\route{POST}{/internal/mcp/list-tools} & Connect to one MCP server and list its tools: the administrator's connection test. \\
\end{xltabular}

The first exception to the third design rule is visible in the table. The public API lets a
browser store only messages of the role \code{user} (\code{GUI_WRITABLE_ROLES} in
\file{chat/api/views.py}); an assistant message, whose tool results the model is shown again
on the next turn, can be written through the internal door alone. A browser can therefore
never plant a forged tool result in a conversation.

Three things protect a door. It is reachable only on the isolated network. The endpoint
forwards no \code{/internal/} path (\cref{sec:arch:network}). And every request must carry a
shared secret.

\begin{codebox}{How the shared secret is checked}
The secret is one value, \code{ASSISTANT_INTERNAL_SECRET}, minted as random hexadecimal by
the update script and given to both containers by the compose file. A caller sends it in the
header \code{X-Internal-Secret}.

\begin{itemize}
  \item \textbf{In the API} the decorator \code{require_internal}
    (\file{assistant/api/internal_views.py}) wraps every view of the door. It compares the
    presented and the expected value as bytes with \code{hmac.compare_digest}, which takes
    the same time wherever the two differ, so the answer's timing tells a guesser nothing
    about how much of a guess was right.
  \item \textbf{In the assistant service} the middleware \code{requireInternal}
    (\file{middleware/internal.js}) stands in front of every \code{/internal} path, before
    the body of the request is parsed. It compares with \code{timingSafeEqual}; a value of
    another length is refused before the comparison, which reveals the length and nothing
    else.
  \item \textbf{Both fail closed.} When the configured secret is empty, every internal call is
    refused; neither side treats \enquote{no secret set} as \enquote{no secret needed}. The
    assistant service also says so in its log when it starts.
\end{itemize}
A refused call is answered 403 with the code \code{INTERNAL_ONLY}.
\end{codebox}

\section{Four requests, followed through the system}\label{sec:arch:flows}

The four flows below are the ones every other chapter builds on. They are drawn as sequences:
time runs downwards, a solid arrow is a step of the main path, a dashed arrow is a request
together with its answer. What happens inside each step --- what a model is given and what is
done with its answer --- belongs to the chapters named in each subsection.

\subsection{A chat turn}

A chat turn is the only flow in which the browser talks to the assistant service
(\cref{fig:arch:chat}). The console sends \route{POST}{/api/assistant/chat} with the user's
token and nothing but the new message: the conversation itself is kept by the API, so the
history a model sees is always the stored one and not what a browser claims it was.

\begin{figure}[H]
\centering
\begin{tikzpicture}[x=1mm, y=1mm]
  \begin{scope}[on background layer]
    \draw[lane] (53,-76) rectangle (150,-106);
  \end{scope}
  \node[lanelabel, anchor=south west] at (60,-106) {repeated while the model calls tools};

  \node[svc] (hc) at (10,0) {console};
  \node[svc] (ha) at (58,0) {assistant service};
  \node[svc] (hd) at (108,0) {API (Django)};
  \node[ext] (hg) at (140,0) {gateway};
  \foreach \x in {10,58,108} \draw[draw=rule, line width=.6pt] (\x,-4) -- (\x,-131);
  \draw[draw=rule, line width=.6pt] (140,-4) -- (140,-78);

  \draw[flow] (10,-13) -- node[lbl, above, align=center] {\code{POST}~\code{/api/assistant/chat}\\ the token, the new message} (58,-13);
  \draw[call] (58,-21) -- node[lbl, above] {\code{/api/auth/me}: who is asking} (108,-21);
  \node[chk, minimum height=5mm, inner ysep=2pt] at (58,-28.5) {turn limit, in memory};
  \draw[call] (58,-37) -- node[lbl, above] {\code{/internal/assistant/budget}} (108,-37);
  \draw[call] (58,-44) -- node[lbl, above] {\code{/internal/assistant/prompt}} (108,-44);
  \draw[call] (58,-54) -- node[lbl, above, align=center] {\code{/api/chat/threads/…}: the thread, its\\ history; the new message is stored} (108,-54);
  \draw[call] (58,-64) -- node[lbl, above, align=center] {\code{/api/companies/…}, \code{/api/reports}:\\ the company and its reports} (108,-64);
  \draw[call] (58,-71) -- node[lbl, above] {\code{/internal/assistant/mcp-servers}} (108,-71);

  \node[ai, text width=15mm, minimum height=24mm] (m) at (140,-91) {the model answers, in at most 24 steps};
  \draw[flow] (58,-82) -- node[lbl, above, pos=.3] {instructions, history, tools} (m.west |- 0,-82);
  \draw[flow] (m.west |- 0,-89) -- node[lbl, above, pos=.72] {a tool call} (58,-89);
  \draw[call] (58,-94) -- node[lbl, above] {\code{/api/*} with the user's token} (108,-94);
  \draw[flow] (58,-99) -- node[lbl, above, pos=.3] {the tool's result} (m.west |- 0,-99);

  \draw[flow] (58,-112) -- node[lbl, above] {the answer, as an event stream} (10,-112);
  \draw[call] (58,-120) -- node[lbl, above] {\code{…/threads/{id}/messages}: the reply} (108,-120);
  \draw[call] (58,-127) -- node[lbl, above] {\code{/internal/assistant/turn-log}} (108,-127);
\end{tikzpicture}
\actorlegend
\caption{A chat turn. The assistant service identifies the caller and loads the context through
  the public API as the user, reads the prompt, the budget and the MCP servers through the
  internal door, and runs the model with tools. Each tool is a request to the public API with
  the user's token. The reply and the usage record go back through the internal door.}\label{fig:arch:chat}
\end{figure}

The route (\file{routes/chat.js}) passes a turn through its gates before any model is called.
The caller's token is resolved against \route{GET}{/api/auth/me}, and the answer is cached
for 15~seconds (\code{IDENTITY_CACHE_MS}); a caller without a valid token is refused. The
turn limiter allows 12 turns a user a minute (\code{RATE_LIMIT_TURNS}). The user's budget for
the day is asked from the API, and a spent budget refuses the turn; if the question itself
fails, the turn is let through rather than every chat be blocked. The active version of the
prompt \prompt{chat.assistant} is read from the API, and without one the chat answers 503:
the assistant is then switched off on purpose.

The service then loads, as the user, the thread, the branch of its history, the company the
thread is about and that company's reports, stores the user's message, and opens the tools of
the enabled MCP servers. It builds the instructions and calls the model on the gateway with
the tools attached. The model may answer or call a tool; a tool call becomes a request to the
public API under the user's token, its result goes back to the model, and this repeats for at
most 24 model steps (\code{MAX_STEPS}), the last of which runs without tools so that the turn
ends in text. The answer streams to the browser as it is written.

When the turn ends the service stores its reply and logs the turn. A turn outlives the
request that started it: if the browser reloads or the tab is closed, the model goes on and
the reply is stored; apart from the turn's own time limits, only the stop button, a newer
turn in the same thread or a shutdown ends it early. On the first turn of a thread, while it
still carries its automatic title, a second, small model call gives the conversation a
title. The tools, the prompt and both model calls are the subject of \cref{ch:chat}.

\subsection{A report}

A report is far too long for a request: it is a job (\cref{fig:arch:report}).
\route{POST}{/api/reports} checks that the caller may see the company, that the corpus is not
empty, that no other report of the company is under way and that the user has budget left;
it then creates the report's row in the state \code{queued}, enqueues a job of the kind
\code{report.analyze} with the report's id as its payload, and answers at once.

\begin{figure}[H]
\centering
\begin{tikzpicture}[x=1mm, y=1mm,
    cell/.style={minimum height=5.6mm, inner ysep=1.5pt, inner xsep=2pt, text width=18mm}]
  \node[svc, text width=17mm] (hc) at (10,0) {console or chat tool};
  \node[svc] (hd) at (44,0) {API (Django)};
  \node[store] (hj) at (70,0) {\code{jobs}};
  \node[svc] (hw) at (100,0) {worker's loop};
  \node[svc, text width=17mm] (ha) at (140,0) {assistant service};
  \foreach \x in {44,70,100} \draw[draw=rule, line width=.6pt] (\x,-5) -- (\x,-128);
  \draw[draw=rule, line width=.6pt] (10,-6) -- (10,-128);
  \draw[draw=rule, line width=.6pt] (140,-6) -- (140,-128);

  \draw[flow] (10,-14) -- node[lbl, above] {\code{POST}~\code{/api/reports}} (44,-14);
  \draw[flow] (44,-21) -- node[lbl, above] {job queued} (70,-21);
  \draw[call] (100,-28) -- node[lbl, above] {claims the job} (70,-28);

  \node[svc, cell] at (89.5,-36) {profile};
  \node[ai, cell] at (110.5,-36) {documents};
  \node[ai, cell] at (89.5,-43) {citations};
  \node[ai, cell] (p1) at (110.5,-43) {applicability};
  \draw[call] (121.2,-39.5) -- node[lbl, above] {\code{/internal/llm/*}} (140,-39.5);

  \draw[flow] (100,-53) -- node[lbl, above, align=center] {job done; the report\\ waits in \code{scope_review}} (70,-53);
  \draw[call] (10,-61) -- node[lbl, above] {polls \code{GET}~\code{/api/reports/{id}}} (44,-61);

  \node[man, text width=17mm] (person) at (10,-73) {a person approves the scope};
  \draw[flow] (person.east) -- node[lbl, above] {\code{POST}~\code{…/scope}} (44,-73);
  \draw[flow] (44,-80) -- node[lbl, above] {queued again} (70,-80);
  \draw[call] (100,-87) -- node[lbl, above] {claims the job} (70,-87);

  \node[ai, cell] at (89.5,-95) {documents};
  \node[ai, cell] at (110.5,-95) {obligations};
  \node[ai, cell] at (89.5,-102) {gap\_analysis};
  \node[ai, cell] at (110.5,-102) {internal};
  \node[ai, cell] at (89.5,-109) {action\_plan};
  \node[ai, cell] at (110.5,-109) {summary};
  \draw[call] (121.2,-102) -- node[lbl, above] {\code{/internal/llm/*}} (140,-102);

  \draw[flow] (100,-118) -- node[lbl, above] {job done; the report is done} (70,-118);
  \draw[call] (10,-125) -- node[lbl, above] {reads it, exports it} (44,-125);
\end{tikzpicture}
\actorlegend
\caption{A report started with the scope review. The job runs twice: the first run ends when
  the acts have their verdicts and leaves the report waiting; a person's approval queues the
  same kind of job again, and the second run assesses the obligations. The boxes on the
  worker's line are the job's steps; a yellow step makes model calls through the assistant
  service. The second run reads the company again, freezes the profile a second time and
  takes the documents step again, which finds the inventories in
  their cache; the step \code{internal} runs only for a report started with the company's own
  requirements. Without the review one run does both parts.}\label{fig:arch:report}
\end{figure}

The worker claims the job and runs its handler, \code{analyze} in \file{analysis/pipeline.py}.
A report is made in two parts. The first decides which acts apply and is cheap: the steps
\code{profile}, \code{documents}, \code{citations} and \code{applicability}. The second
judges every obligation of those acts and is where the money goes: \code{obligations},
\code{gap_analysis}, \code{action_plan} and \code{summary}, with a step \code{internal}
between the gap analysis and the plan when the report was started with the company's own
requirements. Every model call of a step travels from the worker to the assistant service's
door and from there to the gateway.

\begin{humanbox}{The stop between the two parts}
A report started with the review option does not go on by itself. When the first part has
finished, the handler sets the report to \code{scope_review} and returns; the job ends as
done, and nothing runs. The console, which polls \route{GET}{/api/reports/{id}} every three
seconds while a report is active, shows the acts with what the report would do with each. A
person may change the selection and is shown its price (\route{POST}{/api/reports/{id}/scope}
without \code{approved}), and approves it (the same route with \code{approved}). The approval
is stamped on the report with the person and the time, written to the audit log with every
change made, and a second \code{report.analyze} job is enqueued. The handler tells the second
run from the first by that stamp. Without the review option the handler stamps the time
itself, with no person, and goes straight on.
\end{humanbox}

The second run reads what the first left in the database --- the acts and their verdicts and
the citation audit. The company it reads again, as it stands when the second run starts, and
it freezes the profile a second time (\code{freeze_profile}), so that the report shows the
company its findings were judged on: a credential verified or a document added while the
report waited counts. Its \code{documents} step
runs again and finds the documents' inventories in their cache. When the last step has finished the report is \code{done}; a failure sets it
to \code{failed} with a message a user can read, a cancellation to \code{cancelled}. The
finished report is read through \route{GET}{/api/reports/{id}} and exported through
\route{GET}{/api/reports/{id}/export}. The pipeline is the subject of \cref{ch:report}.

\subsection{A work product}

A work product is made the same way, by a job of the kind \code{programme.generate}
(\cref{fig:arch:product}). \route{POST}{/api/companies/{id}/programme} with the kind ---
\code{policy}, \code{governance} or \code{risks} --- checks that the company has readable
documents, that no product of that kind is being generated and that budget is left, creates
the next version of the product in the state \code{queued} and enqueues the job. The chat's
tool \code{startWorkProduct} calls the same route.

\begin{figure}[H]
\centering
\begin{tikzpicture}[x=1mm, y=1mm,
    cell/.style={minimum height=5.6mm, inner ysep=1.5pt, inner xsep=2pt, text width=30mm}]
  \node[svc, text width=17mm] (hc) at (10,0) {console or chat tool};
  \node[svc] (hd) at (48,0) {API (Django)};
  \node[svc] (hw) at (92,0) {worker's loop};
  \node[svc, text width=17mm] (ha) at (140,0) {assistant service};
  \draw[draw=rule, line width=.6pt] (10,-6) -- (10,-72);
  \draw[draw=rule, line width=.6pt] (48,-5) -- (48,-72);
  \draw[draw=rule, line width=.6pt] (92,-5) -- (92,-72);
  \draw[draw=rule, line width=.6pt] (140,-6) -- (140,-72);

  \draw[flow] (10,-14) -- node[lbl, above] {\code{POST}~\code{…/programme}} (48,-14);
  \draw[flow] (48,-21) -- node[lbl, above] {job queued, then claimed} (92,-21);
  \node[ai, cell] at (92,-28) {documents};
  \node[ai, cell] (calls) at (92,-35) {two model calls};
  \node[chk, cell] (refs) at (92,-42) {references};
  \draw[call] (calls.east) -- node[lbl, above] {\code{/internal/llm/generate}} (140,-35);
  \draw[opt] (refs.east) -- node[lbl, below] {the citation judge} (140,-42);
  \draw[flow] (92,-51) -- node[lbl, above] {the product is a draft} (48,-51);
  \draw[call] (10,-58) -- node[lbl, above] {polls \code{GET}~\code{/api/programme/{id}}} (48,-58);
  \node[man, text width=17mm] (person) at (10,-68) {a person edits, approves};
  \draw[flow] (person.east) -- node[lbl, above] {\code{PUT}~\code{…/content}, \code{PATCH}} (48,-68);
\end{tikzpicture}
\actorlegend
\caption{A work product. The job reads the company's documents, asks the model twice and
  puts the answer through the rules of the product's kind; the result is a draft. Only a
  person's request approves it.}\label{fig:arch:product}
\end{figure}

The handler (\code{generate} in \file{programme/run.py}) has three stages. The step
\code{documents} gathers the company's profile and documents and each document's inventory,
from its cache or, when there is none yet, by a model call (\code{ensure_inventory}). The
kind's own two steps make
the two model calls, the second given what the first wrote. The step \code{references}
brings the answer into the kind's shape and checks its rules with code; the provisions the
text names are looked up in the corpus, and whether each regulates what the text says of it
is judged by a model, once. The product is then stored as a \code{draft} with its content and
the result of every rule. It never leaves that state by itself: a person edits it and
approves it through the API (\cref{ch:programme}).

\subsection{A refresh of the legal corpus}

The corpus is kept current by jobs of the kind \code{corpus.ingest}, one per act
(\cref{fig:arch:refresh}). An administrator queues one with
\route{POST}{/api/corpus/documents/{id}/refresh}; once a week the scheduler queues
\code{corpus.refresh_all}, which does nothing but enqueue one \code{corpus.ingest} for every
enabled act whose text is fetched rather than uploaded.

\begin{figure}[H]
\centering
\begin{tikzpicture}[x=1mm, y=1mm,
    cell/.style={minimum height=5.6mm, inner ysep=1.5pt, inner xsep=2pt, text width=34mm}]
  \node[svc, text width=22mm] (hc) at (13,0) {administrator, or the weekly schedule};
  \node[svc] (hw) at (64,0) {worker's loop};
  \node[ext, text width=20mm] (hs) at (108,0) {the source, through the exit};
  \node[svc, text width=17mm] (ha) at (140,0) {assistant service};
  \draw[draw=rule, line width=.6pt] (13,-6) -- (13,-64);
  \draw[draw=rule, line width=.6pt] (64,-5) -- (64,-64);
  \draw[draw=rule, line width=.6pt] (108,-6) -- (108,-64);
  \draw[draw=rule, line width=.6pt] (140,-6) -- (140,-64);

  \draw[flow] (13,-14) -- node[lbl, above] {\code{corpus.ingest} queued} (64,-14);
  \draw[call] (64,-21) -- node[lbl, above] {fetch} (108,-21);
  \node[chk, cell] at (64,-28) {unchanged? then stop here};
  \node[svc, cell] at (64,-35) {extract, parse, chunk};
  \node[ai, cell] (emb) at (64,-42) {embed};
  \node[chk, cell] at (64,-49) {activate the version};
  \draw[call] (emb.east) -- node[lbl, above, pos=.72] {\code{/internal/llm/embed}} (140,-42);
  \draw[flow] (64,-59) -- node[lbl, above] {the act has a new current version} (13,-59);
\end{tikzpicture}
\actorlegend
\caption{A refresh of one act. The worker fetches the source through the exit and stops when
  nothing changed. Otherwise it builds a new version beside the current one and switches to
  it in one transaction.}\label{fig:arch:refresh}
\end{figure}

The handler (\code{ingest} in \file{corpus/ingest.py}) records five steps. \code{fetch}
gets the bytes from the act's source and compares them with the current version: the same
content hash, or the same edition label, ends the job at once with the answer
\enquote{unchanged}, which is why a weekly check of an unchanged corpus costs nothing.
Otherwise a new version row is created and \code{extract} reads the publication's structure
in a child process with limits on memory and time, \code{parse} writes the provisions as
units, \code{chunk} cuts them for retrieval and \code{embed} has the chunks embedded. Only
then is the version made current, together with the record of what it changed against the
one before, in a single transaction; until that moment every reader sees the old version
whole. A run that fails leaves its version marked failed, and a retry starts a fresh one.
For an uploaded standard the \code{extract} step may also ask a model to mark the headings.
All of this is the subject of \cref{ch:corpus}.

\section{Background jobs}\label{sec:arch:jobs}

Everything that takes longer than a request is a \term{job}: one row in a table, picked up
and run by the \term{worker}. There is no message broker and no separate scheduler. The
queue is two tables of the same database that holds everything else, and the code is the
application \code{jobs} of the API: \file{jobs/models.py} (the tables), \file{jobs/queue.py}
(enqueueing), \file{jobs/registry.py} (the kinds), \file{jobs/runner.py} (claiming and
running), \file{jobs/schedule.py} (the periodic kinds) and
\file{jobs/management/commands/run_worker.py} (the loop).

\subsection{The queue: two tables}

A job is a row of \code{jobs} (\cref{tab:arch:jobrow}). Its \term{kind} names the handler
that will run it, its payload is the handler's input, and its \term{subject} says what it is
about --- \code{report:<id>}, \code{corpus:<id>}, \code{document:<id>} --- so that the pages
can list the jobs of one thing and a route can see that one is already under way.
\code{enqueue} (\file{jobs/queue.py}) writes the row; it refuses a kind that has no handler,
so that a mistyped kind fails where it is written and never waits in the table for ever.

\begin{table}[H]
\caption{The columns of \code{jobs} that make it a queue}\label{tab:arch:jobrow}
\begin{tabularx}{\textwidth}{@{}L{37mm}Y@{}}
\toprule
\hd{Column} & \hd{Meaning} \\
\midrule
\code{kind}, \code{payload}, \code{subject} & Which handler, with what input, about what. \\
\code{status} & \code{queued}, \code{running}, \code{done}, \code{failed} or \code{cancelled}; a check constraint allows no other value. \\
\code{priority}, \code{run_after} & The order of claiming, and the earliest moment a job may be claimed. \\
\code{attempts}, \code{max_attempts} & How often the job was claimed, and how often its kind allows. \\
\code{locked_by}, \code{heartbeat_at} & Which worker holds the job, and when that worker last said it was alive. \\
\code{cancel_requested} & Set by the API; read by the handler. \\
\code{progress} & \texttt{\{step, percent, message\}}: what the pages show while the job runs. \\
\code{result}, \code{error} & What the handler returned, or why it failed. \\
\code{created_by} & The user who caused the job; empty for the scheduler's. Model calls of the job are charged to this user where they are charged at all. \\
\code{created_at}, \code{started_at}, \code{finished_at} & The three moments of its life. \\
\bottomrule
\end{tabularx}
\end{table}

A handler divides its work into \term{steps}, and each step is a row of \code{job_steps}:
its ordinal and name, its status (\code{running}, \code{done} or \code{failed} as the runner
writes them), a reference to its input, its output, its duration and its error --- and the
model, the prompt version and the usage of the model calls made in it. The step is the unit
at which a report's cost and provenance are read: \enquote{the gap analysis ran on this
version of this prompt, made this many calls and cost this much}.

\begin{codebox}{What a step records about its model calls}
A handler never counts its own tokens. Every model call is logged with the id of the step it
was made for, and when a step ends --- done or failed --- the runner reads the step's rows
back from \code{llm_calls} in one grouped query (\code{_usage_from_calls} in
\file{jobs/runner.py}) and writes onto the step: the number of calls, the input and output
tokens and the cost, in total and per model; the model that made the most calls; and the
prompt version, when all the calls ran on one. A step that fails half-way therefore still
shows what it had spent. The prompt version is kept as plain text and not as a foreign key,
so the record survives the deletion of the version it names.
\end{codebox}

\subsection{The life of a job}

\Cref{fig:arch:jobstates} shows the states of a job and every way from one to another. The
rest of this subsection goes through them.

\begin{figure}[H]
\centering
\begin{tikzpicture}[x=1mm, y=1mm]
  \node[svc, minimum width=22mm] (q) at (14,0) {queued};
  \node[svc, minimum width=22mm] (r) at (74,0) {running};
  \node[svc, minimum width=22mm] (d) at (136,15) {done};
  \node[svc, minimum width=22mm] (f) at (136,0) {failed};
  \node[svc, minimum width=22mm] (c) at (136,-15) {cancelled};

  \draw[flow] (q) -- node[lbl, above] {claimed: \code{attempts} + 1} (r);
  \draw[flow] (r.north east) -- node[lbl, pos=.55] {the handler returns} (d.west);
  \draw[flow] (r.east) -- node[lbl, pos=.5, align=center] {an error or a dead worker,\\ no attempt left} (f.west);
  \draw[flow] (r.south east) -- node[lbl, pos=.55] {cancel seen at a check} (c.west);
  \draw[flow] (r.north) to[out=150, in=30] (q.north);
  \node[note, anchor=south, align=left] at (44,15) {back to the queue:\\ an error with an attempt left (30 s later)\\ the worker is stopping (the attempt is given back)\\ a heartbeat older than 120 s, with an attempt left};
  \draw[flow] (q.south) |- node[lbl, pos=.75] {cancelled while still queued} ([yshift=-2mm]c.west);
\end{tikzpicture}
\caption{The states of a job. A job is claimed from the queue, and ends done, failed or
  cancelled; three events send a running job back to the queue instead.}\label{fig:arch:jobstates}
\end{figure}

\paragraph{Claiming and locking.} The worker claims one job at a time (\code{claim} in
\file{jobs/runner.py}): among the rows that are \code{queued} and whose \code{run_after} has
passed, the one with the highest \code{priority}, and among equals the oldest. The row is
selected with PostgreSQL's \texttt{FOR UPDATE SKIP LOCKED}, which locks it for the length of
the transaction and makes any other worker asking at the same instant skip it and take the
next. Inside the same transaction the row is set to \code{running}, \code{locked_by} is set
to the worker's id (the container's host name and eight random characters),
\code{heartbeat_at} to the present, and \code{attempts} is raised by one. Two workers can
therefore never run the same job, and running several loops side by side is safe. An index
over status, \code{run_after}, priority and creation time serves the query.

\paragraph{Priority and delay.} Priority is a plain integer, default 0, higher first. The
code uses four values: 5 for indexing an uploaded document, so that a document a user has
just dropped into the chat is searchable before a long corpus job gets its turn; 0 for
everything a user or an administrator starts through the API (and for the ingestion jobs of
the command \code{seed_corpus}); $-1$ for the jobs the bulk commands
\code{extract_obligations}, \code{verify_obligations} and \code{read_sanctions} queue over
the whole corpus (\code{PRIORITY} in each), so that work a user starts goes first; and $-5$
for the scheduler's. \code{run_after} delays
a job: a retry is pushed 30~seconds ahead (\code{BACKOFF_S}), and an indexing job that finds
its user's daily budget spent queues itself again for five minutes after midnight UTC, when
the budget resets.

\paragraph{Attempts.} Each kind is registered with the number of attempts it is allowed
(\cref{tab:arch:kinds}). When a handler raises an error and the job has an attempt left, the
job goes back to the queue with the error text and the delay; otherwise it ends as
\code{failed}. Most kinds allow one attempt; fetching an act (\code{corpus.ingest}) allows two
and indexing a document (\code{documents.index}) three. An error raised while the job's
cancel flag is set ends it as \code{cancelled} instead, whatever attempts are left.

\paragraph{The heartbeat, and a worker that died.} While a handler runs, a second thread
stamps \code{heartbeat_at} every 15~seconds (\code{HEARTBEAT_S}). At the top of every turn of
its loop the worker looks for running jobs whose heartbeat is older than 120~seconds
(\code{STALE_S}): their worker is taken to be dead (\code{recover_stale}). Such a job goes
back to the queue if it has an attempt left. If it has none it is set to \code{failed} with
the message that its worker stopped unexpectedly, its still-open steps are failed with it
--- each with what its model calls had used --- and the kind's \code{on_abandon} hook runs.
The hook exists for the two kinds whose job has a row of its own that mirrors it: it sets
the report, or the work product, to failed as well, so that neither stays \enquote{running}
for ever with no job behind it.

\begin{codebox}{Fencing: no write without the lock}
A worker that was only slow --- paused for minutes, not dead --- may wake up after another
worker has recovered its job. Every write that settles a job --- done, failed, cancelled,
back to the queue --- is therefore conditional on the row still being \code{running} and
still carrying that worker's \code{locked_by} (\code{_fenced} in \file{jobs/runner.py}), and
a progress write on the row still carrying that \code{locked_by}; an update that matches no
row changes nothing. The heartbeat thread notices the same thing first: when its stamp matches no
row it raises a flag, the handler stops at its next check with \code{JobLost}, and nothing
more is written for that attempt.
\end{codebox}

\paragraph{Cancelling.} \route{POST}{/api/jobs/{id}/cancel} (\file{jobs/api/views.py}) does
one of two things. A job still queued is set to \code{cancelled} at once and its kind's
\code{on_abandon} hook settles the report or product behind it. A running job only gets its
\code{cancel_requested} flag; the handler sees it at its next call of
\code{ctx.check_cancelled()}, which handlers make between their expensive phases, raises
\code{JobCancelled}, and the job ends as \code{cancelled} --- not as a failure, and without a
retry. The two writes are conditional and made in that order, so a job that a worker claims
between them still receives the flag. A job that has already finished answers 409. Reports
and work products have cancel routes of their own that do the same to their job. A user may
cancel the jobs they created, an administrator any job; the same rule decides which jobs
\route{GET}{/api/jobs} lists.

\paragraph{Progress.} A handler reports where it is with \texttt{ctx.progress(step, percent, message)}, which writes the three values onto the job's row. The jobs page and the progress
panels of a report and of a work product poll the job and show them.

\subsection{The kinds of job}

A handler is an ordinary function registered under its kind with the decorator
\code{@register} (\file{jobs/registry.py}); \file{jobs/handlers.py} imports every module that
registers one, in the web process and in the worker alike, since the web process must know
the kinds to enqueue them. \Cref{tab:arch:kinds} lists them all.

\begin{xltabular}{\textwidth}{@{}L{37mm}YL{43mm}@{}}
\caption{The registered kinds of job}\label{tab:arch:kinds}\\
\toprule
\hd{Kind, handler, attempts} & \hd{What it does} & \hd{Who can start it} \\
\midrule
\endfirsthead
\toprule
\hd{Kind, handler, attempts} & \hd{What it does} & \hd{Who can start it} \\
\midrule
\endhead
\bottomrule
\endlastfoot
\code{report.analyze}\newline \file{analysis/pipeline.py}\newline 1 attempt, abandon hook & One run of a report: its first part, its second, or both (\cref{fig:arch:report}). & Any user, for a company they may see: \route{POST}{/api/reports} and the approval of the scope; the chat's tools do the same. \\
\code{programme.generate}\newline \file{programme/run.py}\newline 1 attempt, abandon hook & A new version of one work product. & Any user, for a company they may see: \route{POST}{/api/companies/{id}/programme}; the chat's tool. \\
\code{documents.index}\newline \file{companies/indexing.py}\newline 3 attempts & Cuts one uploaded document into passages and has them embedded. Waits for the next day when the uploader's budget is spent. & Any user, by uploading a document (\route{POST}{/api/documents}); the job itself, when it waits. \\
\code{corpus.ingest}\newline \file{corpus/ingest.py}\newline 2 attempts & Fetches one act, or reads an uploaded file, and makes it the act's current version (\cref{fig:arch:refresh}). & An administrator: adding an act with a source, its refresh, an upload. The kind below. The command \code{seed_corpus}. \\
\code{corpus.refresh_all}\newline \file{corpus/ingest.py}\newline 1 attempt & Enqueues one \code{corpus.ingest} for every enabled act that is fetched. & The scheduler, weekly. An administrator: \route{POST}{/api/corpus/refresh-all}. \\
\code{corpus.reembed}\newline \file{corpus/ingest.py}\newline 1 attempt & Embeds every chunk of the current versions again and, unless one act is named, every obligation in both languages and every company document. Needed after the embedding model changes. & An administrator: \route{POST}{/api/corpus/reembed}. \\
\code{corpus.extract_obligations}\newline \file{corpus/obligations.py}\newline 1 attempt & Reads the obligations out of one act's current version. & An administrator, per act. The command \code{extract_obligations}, for every act that is behind. \\
\code{corpus.verify_obligations}\newline \file{corpus/verify.py}\newline 1 attempt & Has the judge compare one act's obligations with the provisions they cite. & An administrator, per act. The command \code{verify_obligations}. \\
\code{corpus.read_sanctions}\newline \file{corpus/sanctions.py}\newline 1 attempt & Reads the sanction provisions of the acts that were never read, one act after another. & The command \code{read_sanctions} only. \\
\code{ops.maintenance}\newline \file{ops/maintenance.py}\newline 1 attempt & The daily purges (\cref{tab:arch:purges}). & The scheduler, daily. \\
\code{ops.noop}\newline \file{jobs/handlers.py}\newline 1 attempt & A smoke test: a number of empty steps, with an optional pause and a forced failure. & The test suite; a person at a shell. \\
\end{xltabular}

No route of the API enqueues an arbitrary kind. A job comes into being only as the side
effect of a route that has checked the caller and the subject, of a management command run on
the host, or of the scheduler. The routes that refresh an act, extract its obligations or have
the judge read them answer the job already queued or running for that act instead of adding
a second.

Not every long operation appears in the table under its own kind: what a report needs is
done inside the report's job. A report whose acts were never read for obligations extracts
them in its own \code{obligations} step, and reads their sanction provisions there too; it
does not wait for another job.

\subsection{The worker and its loop}

The loop is a management command of the API, \code{run_worker}. Each turn it marks a
heartbeat file, lets the scheduler look once a minute, recovers jobs whose worker died,
claims one job and runs it to its end; with nothing to claim it waits two seconds
(\code{--poll}) and looks again. One loop runs one job at a time.

\paragraph{Why the worker is a Docker client.} The loop needs exactly what the API has: the
same code, the same settings and secrets, the same upload folder, the same connection to the
database, the same pinned hosts for fetching acts. The compose file gives it all of that by
running the loop \emph{inside} the API's container. The container \code{legalops-worker}
is only the stock Docker command-line image with one fixed command, \texttt{docker exec legalops-django python3 manage.py run\_worker}, and the host's Docker socket mounted into
it. It has no network at all, and its log is the loop's log. Nothing is built or configured
twice, and a change of the API's code or environment cannot leave the worker behind. The
cost is the socket: access to it amounts to root on the host, which is why that container
has no network and runs nothing else (\cref{ch:security}).

\paragraph{Taking over a running loop.} This arrangement has one awkward property that the
code handles explicitly. A loop started with \texttt{docker exec} belongs to the API's
container, not to the worker's. When the API's container restarts, the loop dies with it,
the \texttt{docker exec} returns, the worker's container exits and its restart policy starts
it --- and with it a new loop --- again. But when only the worker's container restarts, the
old loop keeps running inside the API's container, and a second one is about to be started
beside it. So a new loop first takes over (\code{take_over} in
\file{jobs/management/commands/run_worker.py}): it
reads the process id the previous loop left in a file (\file{/tmp/worker.pid}), checks in
\file{/proc} that this process really is a \code{run_worker} loop and not an unrelated
process that happens to have the same number, sends it the stop signal, and writes its own id
into the file. It does not wait: the old loop stops claiming at once and finishes only the
step it is in.

\paragraph{Stopping gracefully.} The stop signal only raises a flag. The loop stops
claiming, and the handler that is running sees the flag at its next check and raises
\code{WorkerStopping}. The runner then puts the job back in the queue and gives its attempt
back (\code{attempts} is lowered by one), so that being interrupted is not counted as having
failed; the next loop runs the job again. For a report this means the interrupted part
starts over: the handler begins by deleting the rows an earlier attempt of the same part
left behind, so nothing is doubled.

\begin{notebox}{What an update of the system does to a running job}
The two ways a loop can end have different consequences, and the difference matters for the
kinds that allow a single attempt. A loop that is \emph{asked} to stop --- because a new loop
takes over --- hands its job back with the attempt unspent, and the job runs again. A loop
that is \emph{killed} with the API's container, as it is when the whole stack is taken down
for an update, can hand nothing back: its job stays \code{running} until its heartbeat is two
minutes old, and the next loop then recovers it. \code{corpus.ingest} and
\code{documents.index} have attempts left and are queued again. A \code{report.analyze} or a
\code{programme.generate} job has spent its one attempt, so it ends as \code{failed} with the
message that its worker stopped unexpectedly, and its report or work product is set to failed
with it rather than left hanging. What the interrupted run had already paid for is kept in
its steps and in \code{llm_calls}; the report itself has to be started again.
\end{notebox}

\subsection{The scheduler}

Two kinds are periodic, and \file{jobs/schedule.py} holds the whole schedule as a table of
two entries: \code{ops.maintenance} every day and \code{corpus.refresh_all} every seven days.
Once a minute (\code{SCHEDULE_EVERY_S}) the loop calls \code{tick}, which looks up the
newest job of each kind and enqueues a new one, at priority $-5$ and with the subject
\code{schedule:<kind>}, unless that newest job is still queued or running or was created
within the interval. The schedule thus lives in the queue it feeds. There is no cron
container; the scheduler never queues a kind while one of its jobs is still waiting or running; a refresh an administrator starts by
hand counts as the week's run; and after a long outage each kind is simply due and runs
once. A failure of the tick is logged and never stops the loop.

\subsection{The daily maintenance}

\code{ops.maintenance} (\file{ops/maintenance.py}) removes what has no further use and
answers how much of each it removed (\cref{tab:arch:purges}).

\begin{table}[H]
\caption{What the daily maintenance job removes}\label{tab:arch:purges}
\begin{tabularx}{\textwidth}{@{}L{38mm}YL{43mm}@{}}
\toprule
\hd{What} & \hd{Condition} & \hd{Constant} \\
\midrule
Sessions & Expired. & --- \\
Finished jobs, with their steps & Done, failed or cancelled more than 90 days ago --- except a job a report points at, which is kept as long as the report. & \code{JOB_RETENTION_DAYS}~=~90 \\
Usage records & Rows of \code{llm_calls} older than 400 days, with the chat transcripts they carry. & \code{TELEMETRY_RETENTION_DAYS}~=~400 \\
Orphaned upload files & Files in the documents folder that no document row names and that are more than six hours old: what a crash between writing the file and writing the row leaves behind. & \code{ORPHAN_GRACE_HOURS}~=~6 \\
Rate-limit windows & Past their reset time. & --- \\
Unused chat uploads & Documents attached to no company, uploaded more than seven days ago and named in no chat message: a file dropped into the chat and never sent. At most 500 a run. & \code{ORPHAN_UPLOAD_DAYS}~=~7, \code{ORPHAN_UPLOAD_BATCH}~=~500 \\
\bottomrule
\end{tabularx}
\end{table}

The second row is why a report keeps its provenance: the job a report points at, with the
steps that record each prompt version, model and cost, is not purged while the report exists.
A report made with the scope review points at the job of its second part. The job of its first
part is purged like any other, and nothing reads it after the run: the report's statistics,
with the prompt versions, the models and the cost of both parts, are stored on the report when
the run ends (\cref{ch:data}), and the usage records keep the calls for 400 days.

\section{The technology stack and the two modes}\label{sec:arch:stack}

\Cref{tab:arch:stack} gives the versions the repository pins. The Python packages are pinned
exactly in \file{django/Dockerfile} --- the project has no requirements file --- and the
base images by tag. The two Node projects declare minimum versions in their
\file{package.json}; the exact versions installed are those of each \file{package-lock.json}.

\begin{xltabular}{\textwidth}{@{}L{29mm}Y@{}}
\caption{The technology stack with the versions the repository declares}\label{tab:arch:stack}\\
\toprule
\hd{Part} & \hd{Components and versions} \\
\midrule
\endfirsthead
\toprule
\hd{Part} & \hd{Components and versions} \\
\midrule
\endhead
\bottomrule
\endlastfoot
Ingress and egress & Caddy 2.11 (\code{caddy:2.11-alpine}); Caddy with the layer-4 module (\code{vuknf/caddy-l4:2.11}). \\
API and worker's loop & Python 3.13 (\code{python:3.13-slim}); Django 5.2.17; psycopg 3.3.6 and the \code{pgvector} bindings 0.5.0; gunicorn 26.2.0; django-environ 0.14.0; bcrypt 5.0.0; requests 2.34.2. Reading documents: pypdf 6.19.0, cryptography 50.0.2, python-docx 1.2.0, lxml 6.1.3, PyYAML 6.0.3. Exports: WeasyPrint 70.0, openpyxl 3.1.5, Markdown 3.11. \\
Assistant service & Node 24.16.0; Express 5.1; the AI SDK \code{ai} 7.0 with its provider packages \code{@ai-sdk/openai} 4.0, \code{@ai-sdk/openai-compatible} 3.0 and \code{@ai-sdk/anthropic} 4.0, and \code{@ai-sdk/mcp} 2.0; zod 4.1. \\
Console and public site & React 19.2; Vite 7.3; MUI 7.3 with its data grid 8.27; Tailwind CSS 4.2; assistant-ui 0.15; TanStack Query 5; React Router 7.13; axios 1.13. Tested with vitest 5 and msw 2; in production the build is served by Caddy 2.11. \\
Database & PostgreSQL 17.11 (\code{postgres:17.11-alpine3.24}) with \code{pgvector} 0.8.6 compiled into the image. \\
Worker's container & The Docker command-line client (\code{docker:27-cli}). \\
Database tool & DbGate 7.3.3. \\
\end{xltabular}

The compose file is written to run in one of two modes, and the lines that differ carry a
comment, \texttt{\# Dev} or \texttt{\# Prod}. Switching mode means commenting out the lines of
one tag, uncommenting those of the other and running the update script
(\cref{tab:arch:modes}). The file in the repository stands in dev mode.

\begin{table}[H]
\caption{The two modes of \file{docker-compose.yml}}\label{tab:arch:modes}
\begin{tabularx}{\textwidth}{@{}L{27mm}YY@{}}
\toprule
\hd{Service} & \hd{Dev mode} & \hd{Production mode} \\
\midrule
Console and public site & \file{vite/Dockerfile.dev}: the Vite development server on the source folder, mounted from the host; a saved file reloads the page. & \file{vite/Dockerfile}: the application is built once, from the locked dependencies, and Caddy serves the result as static files. \\
Assistant service & \file{assistant/Dockerfile.dev}: the source mounted from the host, restarted on change by nodemon. & \file{assistant/Dockerfile}: the locked production dependencies and the source baked in; \texttt{node server.js}. \\
API & \code{DJANGO_RUNSERVER} is set: Django's development server, on the source mounted from the host. & The variable is unset: gunicorn from the baked image, by default three processes of eight threads each (\file{django/gunicorn.conf.py}). \\
Networks & The console's and the assistant service's containers also join the external network, to install their packages. & Both are on the isolated network only. \\
\bottomrule
\end{tabularx}
\end{table}

Two consequences of dev mode belong in an architecture chapter because they weaken its two
main properties. A container on the external network can reach the internet directly, so in
dev mode the assistant service is not confined to the exit's allowlist. And the worker's
loop does not reload when a source file changes: it runs the code it started with until the
worker's container is restarted, which makes the new loop take over. In production mode a
change of code reaches any container only through the update script (\cref{ch:ops}).

\section{Architectural decisions}\label{sec:arch:decisions}

The repository records six decisions that shaped the architecture
(\file{docs/markdown/en/architecture.md}, section~10). They are given here with their reasons.

\begin{description}
  \item[Django for the API, not FastAPI.] The record gives this as a decision taken at the
    outset, with no technical argument for it. What the
    code shows is how little of Django is used: the ORM, the migrations, the request cycle
    and the management commands, with none of the bundled applications.

  \item[Chat first; pages for review and administration.] A user is meant to work by
    talking: describe the company, drop its documents, ask for the analysis. The pages of
    the console exist to review what was produced and to administer the system. The chat's
    tools cover the work on a company and nothing else --- they can read the corpus but not
    change it, and cannot reach accounts, the prompts or the MCP settings. This bounds what an instruction hidden in a
    document could make the chat do (\cref{ch:security}).

  \item[The system drafts and checks; a person approves.] The scope of a report, an
    obligation of the register, the override of a finding, a work product: each is approved
    by a named person and audited, and nothing is marked approved by the system. Scores,
    bands and identifiers are computed by code and never taken from a model. The stop for
    the scope review in \cref{fig:arch:report} and the draft state in
    \cref{fig:arch:product} are this decision in the architecture; \cref{ch:control} gathers
    all such points.

  \item[Obligations belong to the corpus, and are reviewed.] The obligations of an act are
    extracted once per version and reused by every later report of every company, instead of
    being read anew for each report. This makes a report cheaper and two reports
    comparable, and it gives the review of an obligation a lasting effect. A report records
    how many of its findings rest on reviewed obligations. How the register was built and
    reviewed is the subject of \cref{ch:obligations}.

  \item[Official machine sources for the law.] The texts come from CELLAR and from e-tar
    rather than from EUR-Lex and e-seimas, the sites a person would read, because those two
    refuse unattended clients. The egress allowlist of \cref{tab:arch:egress} is the
    architectural trace of this decision.

  \item[One model provider: the faculty's gateway.] Every model call goes to the gateway at
    \code{ai.knf.vu.lt}. The record states plainly what leaves the machine as a result: the
    company's documents are sent to the gateway to be embedded and in the steps of a report,
    and whatever a chat reads travels with its turns. Everything else --- the database, the
    files, the parsing, the checks, the exports --- runs on the host. The assistant service
    can be configured for other providers, but the allowlist would have to be opened for
    them (\cref{ch:models}).
\end{description}
